\Page{18}{\textenglish{04 / FUTURE}}{המקצועות שנפתחים מחדש}{44}{51}

{\fontsize{13.2}{18.8}\selectfont

היום אפשר לבנות דוגמה וירטואלית, להפיק גרסאות לפריים בלי להחזיר את כל הצוות לאתר, ולהשוות כמה תחזיות לפני הזמנה. בכל אחד מן המצבים האלה קיים חומר שמישהי יכולה לבחון. תדמיתנית יכולה לסמן מידות שמצריכות תיקון, צלמת יכולה לזהות שהאור בגרסה מחוללת כבר אינו נתפס טבעי, וקניינית יכולה להסביר שדפוס מסוים נובע ממחסור במלאי. עבודה עם מודל נותנת להן תוצר חדש לבדיקה. הידע שלהן קובע אם התיקון הבא יחזיר אותו לכוונה המקורית. \SourceRef{2}\SourceRef{5}\SourceRef{7}\SourceRef{8}

\par}

\Photo{hands-at-work.jpg}{61mm}{ידיים עובדות ליד מכונת תפירה. \TailorCredit}

{\fontsize{12.6}{17.7}\selectfont

\StoryColumns{דמיינו תרחיש. סטודיו מחבר מעצבת, תדמיתנית ואשת נתונים מוקדם יותר בתהליך. יחד הן יכולות לבדוק שלוש משפחות של צלליות לפני הזמנת דוגמאות, ואז להקדיש עבודה פיזית לרעיון הקשה יותר לביצוע. זה מסלול אפשרי. במסלול אחר, מנהלת רואה כמה גרסאות אפשר להפיק ביום ומוסיפה עוד משימות לבדיקת כל גרסה. אותן יכולות יכולות להרחיב את מרחב המשחק או לקצר את הזמן שמוקדש לכל החלטה.}{מחקר ארגון העבודה הבינלאומי מ־2025 מבחין בין חשיפה של משימות לבינה מלאכותית מחוללת לבין אובדן משרות. הוא עוזר לנסח שאלה קונקרטית לתעשייה: מי מבצעת עכשיו את העבודה שנחסכה במקום אחר? אם פחות שעות מוקדשות לתכנון לוקיישנים אבל יותר זמן נדרש לבדיקת דימויים, חשוב לעקוב אחר התקציב, ההכשרה וחלוקת הקרדיט. השינוי הזה לא אומר מראש כמה מעצבות, צלמים או אנשי הפקה יידרשו בעונה הבאה. \SourceRef{13}}

\par}

{\fontsize{12.6}{17.7}\selectfont

מבחן מעניין לסטודיו אחרי עונה יהיה לבדוק איזה רעיון הגיע לדוגמה שבעבר היה נזנח, מי קיבלה זמן לבדוק אותו, וכמה תיקונים נדרשו כדי לאשר אותו. כך אפשר לראות אם יכולת ייצור מהירה של גרסאות נהפכה ליכולת יצירתית. היא תהיה בעלת ערך כשמישהי השתמשה בזמן שהתפנה כדי לעשות דבר שהיה קשה להצדיק קודם.

\par}
